\section{Aufgabenstellung und Anforderungen}
\label{sec:aufgabenstellung}

\subsection{Ausgangslage}
\label{sec:ausgangslage}

Die Zelle verfügt bereits über zwei getrennte, unabhängig voneinander
entstandene Anbindungen an externe Rechner (vgl.\ Abschnitt~\ref{sec:sdt}).
Eine MATLAB-seitige Bahnvorgabe/-auslesung über \ac{kvp} (Friesen) sowie eine
Sensorintegration über \ac{ethercat} (Hübel). Beide Anbindungen laufen aktuell
getrennt voneinander und ohne harte Echtzeitgarantie. Laut Aufgabenstellung
ist das erklärte Ziel von Thema~2, diese beiden Stränge in eine gemeinsame,
deterministisch getaktete QUARC/Simulink-Struktur zu überführen und um eine
kraftgeregelte Anwendung zu erweitern.

Im weiteren Projektverlauf wurde anstelle des ursprünglich vorgesehenen
Kraftsensors ein Ultraschallsensor verwendet. Die
Regelungsapplikation wurde dementsprechend als Abstandsregelung umgesetzt.

\subsection{Anforderungen}
\label{sec:anforderungen}

Aus der ursprünglichen Aufgabenstellung ergeben sich vier
Kernanforderungen:

\begin{enumerate}
    \item \textbf{Analyse und Portierung KVP~$\rightarrow$~QUARC:} Die
    bestehende MATLAB/KVP-Anbindung ist zu analysieren und auf eine
    QUARC-basierte Simulink-Struktur mit deterministischem Abtastverhalten
    zu portieren.
    \item \textbf{Kraftsensor-Integration:} Ein Kraftsensor ist sowohl
    hardwareseitig (Anschluss an die Analogeingänge der QUARC-I/O-Karte) als
    auch softwareseitig in die neue Struktur zu integrieren.
    \item \textbf{Erfassung und Analyse der Ende-zu-Ende-Systemlatenzen:} Die
    Latenzen des Gesamtsystems sind zu messen und zu analysieren -- sowohl
    für den bestehenden MATLAB/KVP-Pfad als Baseline (siehe
    Abschnitt~\ref{sec:baseline}) als auch für die neue QUARC-Anbindung.
    \item \textbf{Regelungsapplikation:} Auf Basis der deterministischen
    Anbindung und der Kraftmessung ist eine reaktive, totzeitkompensierte
    Regelungsapplikation umzusetzen.
\end{enumerate}
